\chapter{Grids: Convolutional Networks}
\label{ch:grids}

\providecommand{\conv}{*}
\providecommand{\Ft}{\mathcal{F}}
\providecommand{\sinc}{\operatorname{sinc}}
\providecommand{\tr}{\operatorname{tr}}
\providecommand{\diag}{\operatorname{diag}}
\providecommand{\supp}{\operatorname{supp}}

Convolutional neural networks (CNNs) are the paradigmatic case of equivariant
neural architectures, where the underlying mathematical structure dictates both
their capabilities and their limitations. This chapter develops the mathematical
foundations of convolutional equivariance: from its historical and empirical
origins to its rigorous formalization, and the spectral perspective that
emerges when convolution is read through the lens of harmonic analysis. We
start from classical CNNs, whose success in computer vision is explained
mathematically by translation equivariance, but whose restriction to this single
group of symmetries is also their fundamental limitation. Once these limitations
are made precise and the connection with harmonic analysis is established, the
Fourier perspective both explains what convolutional filters do and foreshadows
the spectral methods on groups, graphs, and manifolds developed in later
chapters. This progression from an initial restriction towards full generality
is exactly the one predicted by the unifying framework of
\Cref{ch:foundations,ch:priors}: CNNs were identified there as instances of the
message-passing framework over regular grids, and this chapter develops in depth
the mathematical foundations of that connection.

\section{Classical convolutional neural networks}
\label{sec:grids:classical}

CNNs arise in the context of automatic pattern recognition in images. Their
origin dates to the late 1980s and early 1990s, when Yann LeCun demonstrated
their potential for handwritten digit recognition~\citep{lecun1989}.

The idea of convolutional networks is born from the techniques used to detect
patterns in images. The technique was relatively simple: a convolution with a
kernel designed to detect a specific pattern in the image --- for example, a
kernel that detects horizontal, vertical, or diagonal edges. Mathematically,
this translates into the following operation, where $I$ is the image and $K$ is
the kernel:
\begin{equation}
  (I \conv K)(x,y) = \sum_{u=-k}^{k}\sum_{v=-k}^{k} I(x+u,\, y+v)\, K(u,v).
  \label{eq:grids:convolution2d}
\end{equation}
Strictly speaking, this operation is a \emph{cross-correlation}, not a
convolution, since the mathematical convolution uses $I(x-u,\,y-v)\,K(u,v)$
(see \Cref{def:grids:discrete-conv}). In the deep-learning literature, however,
it is conventional to call it ``convolution'': because the kernel weights are
learned, the difference between convolution and cross-correlation is immaterial
--- the learned kernel is simply reflected. We follow this standard convention
throughout.

A classical example illustrates the kind of kernels that were hand-designed for
image processing.

\begin{example}[The Sobel kernels]
\label{ex:grids:sobel}
The \textbf{Sobel kernels} detect horizontal and vertical edges through
approximations of the partial derivatives.

\textbf{Horizontal Sobel} (detects vertical edges):
\[
K_x = \begin{pmatrix} -1 & 0 & 1 \\ -2 & 0 & 2 \\ -1 & 0 & 1 \end{pmatrix}.
\]

\textbf{Vertical Sobel} (detects horizontal edges):
\[
K_y = \begin{pmatrix} -1 & -2 & -1 \\ 0 & 0 & 0 \\ 1 & 2 & 1 \end{pmatrix}.
\]
These kernels approximate the partial derivatives $\partial I/\partial x$ and
$\partial I/\partial y$ respectively, where $I(x,y)$ denotes the image
intensity. The gradient magnitude is computed as
\[
  \norm{\nabla I} \approx \sqrt{(I \conv K_x)^2 + (I \conv K_y)^2}.
\]
\end{example}

Kernels such as the Sobel ones exhibit revealing symmetry properties: edge
detectors are \emph{antisymmetric} along the detection direction and sum to zero
(removing constant components), whereas smoothing kernels are \emph{symmetric}.
These symmetries directly reflect the geometric transformations that each kernel
detects.

The difficulty with hand-designed kernels was finding the specific one for each
pattern to be detected; moreover, such patterns had to be simple. LeCun's
proposed solution was to use a neural network to \emph{learn} the kernels. This
idea can be taken one step further: rather than learning kernels only for simple
patterns, one can detect more complex patterns by applying kernels to the
outputs of other kernels. That is, a series of kernels is learned in the first
layer of the network, and one keeps stacking layers of kernels on top of the
previous ones to detect ever more complex patterns. The early layers learn
simple strokes such as horizontal and vertical lines, but each subsequent layer
learns more abstract concepts --- such as cubes or rectangles --- and so on,
until the network is able to recognize complex objects such as a cat or a dog.

The explosion of convolutional networks and deep learning took place in 2012,
when AlexNet won the ImageNet competition by a significant
margin~\citep{krizhevsky2012alexnet}. The three factors necessary for the
success of deep learning came together: a large dataset, a neural-network
architecture using deep-learning techniques, and the processing power of
graphics processing units (GPUs).

\begin{figure}[t]
\centering
\begin{tikzpicture}[
    layerconv/.style={draw=gdlblue!60, fill=gdlblue!12, minimum height=1.7cm,
      rounded corners=1pt, font=\scriptsize, align=center, inner sep=2pt},
    layerfc/.style={draw=gdlgreen!60, fill=gdlgreen!12, minimum width=0.45cm,
      minimum height=1.7cm, rounded corners=1pt, font=\scriptsize, align=center},
    lbl/.style={font=\scriptsize, text=gdlgray!75!black, align=center},
    arrow/.style={-{Stealth[length=4pt]}, gdlgray!70, thick},
]
\node[layerconv, minimum width=1.1cm] (c1) at (0,0) {Conv\\$11{\times}11$};
\node[layerconv, minimum width=0.95cm, right=0.35cm of c1] (c2) {Conv\\$5{\times}5$};
\node[layerconv, minimum width=0.8cm, right=0.35cm of c2] (c3) {Conv\\$3{\times}3$};
\node[layerconv, minimum width=0.7cm, right=0.3cm of c3] (c4) {Conv\\$3{\times}3$};
\node[layerconv, minimum width=0.6cm, right=0.3cm of c4] (c5) {Conv\\$3{\times}3$};
\node[layerfc, right=0.6cm of c5] (f1) {FC};
\node[layerfc, right=0.25cm of f1] (f2) {FC};
\node[layerfc, fill=gdlorange!18, draw=gdlorange!60, right=0.25cm of f2] (f3) {FC};
\node[lbl, below=0.15cm of c1] {input\\$224{\times}224$};
\node[lbl, below=0.15cm of c3] {5 convolutional layers\\(with max pooling)};
\node[lbl, below=0.15cm of f2] {3 fully connected\\layers};
\node[lbl, right=0.2cm of f3, anchor=west] {1000-way\\softmax};
\foreach \a/\b in {c1/c2,c2/c3,c3/c4,c4/c5,c5/f1,f1/f2,f2/f3}
  \draw[arrow] (\a) -- (\b);
\end{tikzpicture}
\caption[The AlexNet architecture]{The AlexNet architecture: a convolutional
network with 8 layers (5 convolutional and 3 fully connected) that reached a
top-5 error of $15.3\%$ on ImageNet ILSVRC-2012, a significant improvement over
the $26.2\%$ of the second-best model.}
\label{fig:grids:alexnet}
\end{figure}

From the mathematical viewpoint, the central architectural principle of CNNs is
\emph{weight sharing}: the same kernels are applied at every spatial position,
which is precisely the realization of translation equivariance as an inductive
bias. The local connectivity of each kernel, combined with the hierarchical
composition of layers (from simple edges to abstract representations), allows the
network to progressively build representations of increasing complexity.
However, the restriction to the group of translations is also the fundamental
limitation of these architectures, motivating the geometric extensions developed
in later chapters.

Convolutions can be applied not only to images but extrapolated to any
dimension. For time series we apply one-dimensional convolutions, and for 3D
volumes we apply three-dimensional convolutions. The formulation is analogous:
\begin{equation}
  (I \conv K)(x) = \sum_{u=-k}^{k} I(x+u)\, K(u),
  \label{eq:grids:convolution1d}
\end{equation}
\begin{equation}
  (I \conv K)(x,y,z) = \sum_{u=-k}^{k}\sum_{v=-k}^{k}\sum_{w=-k}^{k}
    I(x+u,\, y+v,\, z+w)\, K(u,v,w).
  \label{eq:grids:convolution3d}
\end{equation}

\subsection{Mathematical foundations and limitations of classical CNNs}
\label{sec:grids:foundations-limits}

\paragraph{Convolution as an equivariant operation.}
The fundamental property that makes CNNs successful is their equivariance to
translations: when the input is translated, the output is translated in the same
way. This property, which we explore through a simple example, is the basis of
their efficiency in image processing.

\begin{example}[Translation equivariance in one dimension]
\label{ex:grids:equiv-1d}
Consider a one-dimensional signal $f \colon \Z \to \R$ and a kernel
$K = [-1, 0, 1]$ (edge detector). The translation action by $\tau \in \Z$ is
given by $(T_\tau f)(n) = f(n - \tau)$.

Let $f = [1, 2, 1, 0, 0]$ and $\tau = 1$. Using
\Cref{eq:grids:convolution1d} with zero-extension outside the domain,
\begin{align}
T_1 f &= [0, 1, 2, 1, 0] && \text{(shifted signal)}, \\
f \conv K &= [2, 0, -2, -1, 0] && \text{(original convolution)}, \\
(T_1 f) \conv K &= [1, 2, 0, -2, -1] && \text{(convolution of shifted signal)}, \\
T_1(f \conv K) &= [1, 2, 0, -2, -1] && \text{(shift of convolution)}.
\end{align}
We observe that $(T_1 f) \conv K = T_1(f \conv K)$, verifying equivariance.
\end{example}

This equivariance property is fundamental for understanding both the strengths
and the limitations of CNNs. The formal notions of equivariance and their
generalization to other transformation groups are developed in detail in
\Cref{ch:priors}.

\begin{figure}[t]
\centering
\begin{tikzpicture}[
    cell/.style={minimum size=0.62cm, inner sep=0pt, outer sep=0pt,
                 draw=gdlgray!50, line width=0.3pt},
    inputcell/.style={cell, fill=gdlblue!15},
    kernelhighlight/.style={cell, fill=gdlorange!25, draw=gdlorange!60, line width=0.5pt},
    kernelcell/.style={cell, fill=gdlorange!20, draw=gdlorange!60, line width=0.5pt},
    outputcell/.style={cell, fill=gdlgreen!15, draw=gdlgray!50},
    outputhighlight/.style={cell, fill=gdlgreen!30, draw=gdlgreen!60, line width=0.6pt},
    lbl/.style={font=\small\bfseries},
]
\def\cs{0.62}
\def\ix{0}\def\iy{0}
\foreach \r [count=\ri from 0] in {{1,0,2,1,3},{2,1,0,3,1},{0,3,1,2,0},{1,2,3,0,1},{3,0,1,2,2}} {
  \foreach \val [count=\ci from 0] in \r {
    \pgfmathsetmacro{\px}{\ix + \ci*\cs + \cs/2}
    \pgfmathsetmacro{\py}{\iy + (4-\ri)*\cs + \cs/2}
    \node[inputcell] at (\px,\py) {\scriptsize\color{gdlblue!70!black}\val};
  }}
\foreach \r [count=\ri from 0] in {{1,0,2},{2,1,0},{0,3,1}} {
  \foreach \val [count=\ci from 0] in \r {
    \pgfmathsetmacro{\px}{\ix + \ci*\cs + \cs/2}
    \pgfmathsetmacro{\py}{\iy + (4-\ri)*\cs + \cs/2}
    \node[kernelhighlight] at (\px,\py) {\scriptsize\color{gdlblue!70!black}\val};
  }}
\draw[gdlorange!80, line width=1.3pt, rounded corners=1pt] (\ix,\iy+2*\cs) rectangle ++(3*\cs,3*\cs);
\node[lbl, text=gdlblue!70!black] at (\ix+2.5*\cs,\iy-0.5*\cs) {Input $\mathbf{X}$};
\pgfmathsetmacro{\starx}{\ix + 5*\cs + 0.5}
\pgfmathsetmacro{\stary}{\iy + 2.5*\cs}
\node[font=\large, text=gdlgray!80!black] at (\starx,\stary) {$\star$};
\pgfmathsetmacro{\kx}{\starx + 0.5}
\pgfmathsetmacro{\ky}{\stary - 1.5*\cs}
\foreach \r [count=\ri from 0] in {{1,0,-1},{1,0,-1},{1,0,-1}} {
  \foreach \val [count=\ci from 0] in \r {
    \pgfmathsetmacro{\px}{\kx + \ci*\cs + \cs/2}
    \pgfmathsetmacro{\py}{\ky + (2-\ri)*\cs + \cs/2}
    \node[kernelcell] at (\px,\py) {\scriptsize\color{gdlorange!70!black}\val};
  }}
\draw[gdlorange!80, line width=1.1pt, rounded corners=1pt] (\kx,\ky) rectangle ++(3*\cs,3*\cs);
\node[lbl, text=gdlorange!70!black] at (\kx+1.5*\cs,\ky-0.5*\cs) {Kernel $\mathbf{W}$};
\pgfmathsetmacro{\eqx}{\kx + 3*\cs + 0.5}
\node[font=\large, text=gdlgray!80!black] at (\eqx,\stary) {$=$};
\pgfmathsetmacro{\ox}{\eqx + 0.5}
\pgfmathsetmacro{\oy}{\ky}
\foreach \r [count=\ri from 0] in {{0,-2,-1},{-1,1,2},{-1,1,2}} {
  \foreach \val [count=\ci from 0] in \r {
    \pgfmathsetmacro{\px}{\ox + \ci*\cs + \cs/2}
    \pgfmathsetmacro{\py}{\oy + (2-\ri)*\cs + \cs/2}
    \ifnum\ri=0 \ifnum\ci=0
      \node[outputhighlight] at (\px,\py) {\scriptsize\bfseries\color{gdlgreen!70!black}\val};
    \else
      \node[outputcell] at (\px,\py) {\scriptsize\color{gdlgreen!70!black}\val};
    \fi \else
      \node[outputcell] at (\px,\py) {\scriptsize\color{gdlgreen!70!black}\val};
    \fi
  }}
\draw[gdlgreen!60, line width=0.8pt, rounded corners=1pt] (\ox,\oy) rectangle ++(3*\cs,3*\cs);
\node[lbl, text=gdlgreen!70!black] at (\ox+1.5*\cs,\oy-0.5*\cs) {Output $(\mathbf{X} \star \mathbf{W})$};
\pgfmathsetmacro{\arcsrcx}{\ix + 1.5*\cs}
\pgfmathsetmacro{\arcsrcy}{\iy + 5*\cs + 0.08}
\pgfmathsetmacro{\arctgtx}{\ox + 0.5*\cs}
\pgfmathsetmacro{\arctgty}{\oy + 3*\cs + 0.08}
\draw[gdlpurple!55, dashed, line width=0.7pt, -stealth, shorten >=1pt]
  (\arcsrcx,\arcsrcy) to[out=40,in=140] (\arctgtx,\arctgty);
\end{tikzpicture}
\caption[The 2D convolution operation]{The two-dimensional convolution
operation: a $3 \times 3$ kernel $\mathbf{W}$ slides over the input image
$\mathbf{X}$, computing the local inner product at each position to produce the
output activation map $(\mathbf{X} \star \mathbf{W})$.}
\label{fig:grids:conv}
\end{figure}

\begin{definition}[Discrete convolution]
\label{def:grids:discrete-conv}
For functions $f, g \colon \Z^d \to \R$, the \emph{discrete convolution} is
\[
  (f \conv g)(x) = \sum_{a \in \Z^d} f(a)\, g(x - a),
\]
where the sum ranges over all points $a$ of the domain $\Z^d$. In practice, if
$f$ has finite support (as for a digital image), the sum reduces to the indices
where $f(a) \neq 0$.
\end{definition}

\begin{definition}[Continuous convolution]
\label{def:grids:continuous-conv}
For functions $f, g \colon \R^d \to \R$, the \emph{continuous convolution} is
\[
  (f \conv g)(x) = \int_{\R^d} f(a)\, g(x - a)\, \mathrm{d}a,
\]
where the integral is taken with respect to the Lebesgue measure over all of
$\R^d$, and exists when $f, g \in L^1(\R^d)$ or under compact-support
conditions.
\end{definition}

The discrete convolution can be seen as a discretization of the continuous one,
where the integral is approximated by a Riemann sum. In the context of digital
image processing we work with the discrete version, while analogue signals and
theoretical analysis use the continuous version.

\begin{proposition}[Algebraic properties of convolution]
\label{prop:grids:algebra}
The discrete convolution satisfies the following properties for finitely
supported functions $f, g, h \colon \Z^d \to \R$:
\begin{enumerate}
  \item \textbf{Commutativity:} $f \conv g = g \conv f$.
  \item \textbf{Associativity:} $(f \conv g) \conv h = f \conv (g \conv h)$.
  \item \textbf{Distributivity:} $f \conv (g + h) = f \conv g + f \conv h$.
  \item \textbf{Identity:} $f \conv \delta = f$, where $\delta$ is the Kronecker
    delta.
\end{enumerate}
\end{proposition}

\begin{proof}
\textbf{(1) Commutativity.} Through the change of variable $b = x - a$,
\[
  (f \conv g)(x) = \sum_{a \in \Z^d} f(a)\,g(x - a)
   = \sum_{b \in \Z^d} f(x - b)\,g(b) = (g \conv f)(x).
\]

\textbf{(2) Associativity.} By direct application of the definition and
exchanging the order of summation (discrete Fubini),
\begin{align}
((f \conv g) \conv h)(x)
  &= \sum_{b} (f \conv g)(b)\,h(x - b)
   = \sum_{b} \Big[\sum_{a} f(a)\,g(b - a)\Big] h(x - b) \nonumber \\
  &= \sum_{a} f(a) \sum_{b} g(b - a)\,h(x - b).
\end{align}
With the change $c = b - a$,
\[
  = \sum_{a} f(a) \sum_{c} g(c)\,h(x - a - c)
   = \sum_{a} f(a)\,(g \conv h)(x - a) = (f \conv (g \conv h))(x).
\]

\textbf{(3) Distributivity.} By linearity of the sum,
\[
  (f \conv (g + h))(x) = \sum_{a} f(a)\,(g + h)(x - a)
   = (f \conv g)(x) + (f \conv h)(x).
\]

\textbf{(4) Identity.} The Kronecker delta satisfies $\delta(a) = 1$ if $a = 0$
and $\delta(a) = 0$ otherwise, so $(f \conv \delta)(x) = \sum_a f(a)\,\delta(x-a)
= f(x)$, since $\delta(x - a) \neq 0$ only when $a = x$.
\end{proof}

These properties have direct architectural consequences. Associativity means
that stacking two convolutional layers (ignoring nonlinearities) is itself a
convolution, so depth composes cleanly; commutativity and the identity make the
algebra of filters as well behaved as ordinary multiplication. The central
result, however, is the link to symmetry.

\begin{theorem}[Translation equivariance of convolution]
\label{thm:grids:equivariance}
Convolution is equivariant to the group of translations $T = \Z^d$:
\[
  T_b(f \conv g) = (T_b f) \conv g = f \conv (T_b g),
\]
where $(T_b f)(x) = f(x-b)$ is translation by $b$.
\end{theorem}

\begin{proof}
For the first equality, with the index change $a' = a - b$,
\begin{align}
\big[(T_b f) \conv g\big](x)
  &= \sum_a (T_b f)(a)\, g(x - a) = \sum_a f(a - b)\, g(x - a) \nonumber \\
  &= \sum_{a'} f(a')\, g(x - a' - b) = (f \conv g)(x - b)
   = \big[T_b(f \conv g)\big](x).
\end{align}
For the second equality,
\begin{align}
\big[T_b(f \conv g)\big](x)
  &= (f \conv g)(x - b) = \sum_a f(a)\, g\big((x - b) - a\big) \nonumber \\
  &= \sum_a f(a)\, g\big((x - a) - b\big) = \sum_a f(a)\,(T_b g)(x - a)
   = \big[f \conv (T_b g)\big](x).
\end{align}
\end{proof}

Translation equivariance implies that the kernel weights must be shared
spatially, which drastically reduces the number of parameters compared with a
fully connected network. This fundamental property is generalized to broader
groups in \Cref{ch:groups}.

\paragraph{The converse: convolution is forced by equivariance.}
What makes convolution \emph{the} grid layer rather than merely \emph{a} grid
layer is that any linear map between signals on $\Z^d$ that commutes with all
translations is necessarily a convolution. A linear translation-equivariant
operator $\Phi$ is determined by its response to the unit impulse $\delta$; call
that response $w = \Phi\delta$. Any signal decomposes as
$f = \sum_a f(a)\, T_a \delta$, and by linearity and equivariance
$\Phi f = \sum_a f(a)\, T_a w = f \conv w$. Thus equivariance to translation does
not merely permit convolution --- it forces it. This is the blueprint of
\Cref{ch:priors} realized in its purest form.

\paragraph{Limitations of classical CNNs.}
The strength of the CNN inductive bias is also its boundary.

\begin{example}[Lack of rotational equivariance]
\label{ex:grids:no-rotation}
Consider an image of the digit ``6'' and its $180^\circ$ rotation, which looks
like a ``9''. A standard CNN must learn these transformations from data through
data augmentation, rather than incorporating them structurally. If $R_{90}$
denotes a rotation by $90^\circ$, then
\[
  f\big(R_{90}(\text{image})\big) \neq R_{90}\big(f(\text{image})\big).
\]
This illustrates how standard CNNs are equivariant only to translations, not to
other geometric transformations.
\end{example}

Although CNNs have revolutionized computer vision, their design exhibits
fundamental limitations that restrict their applicability. These limitations,
which arise directly from their restriction to translation equivariance, can be
classified into four main categories:
\begin{enumerate}
  \item \textbf{Symmetry limitations.} They are equivariant only to discrete
    translations $\Z^d$, requiring data augmentation for transformations such as
    rotations, reflections, or scalings. This is inefficient both
    computationally and statistically.
  \item \textbf{Topological limitations.} Designed for regular Euclidean grids,
    they are inadequate for data with non-trivial topologies such as graphs,
    point clouds, Riemannian manifolds, or simplicial complexes.
  \item \textbf{Locality limitations.} Their reliance on local structures means
    that global relationships require many layers to be captured, limiting
    efficiency in domains where long-range dependencies are critical.
  \item \textbf{Receptive-field limitations.} The linear growth of the receptive
    field with depth requires deep architectures to capture global patterns,
    introducing vanishing-gradient problems and computational inefficiency.
\end{enumerate}

These systematic limitations have motivated the geometric extensions that form
the core of the remaining Gs. Each extension addresses specific subsets of these
limitations through the controlled generalization of the preserved symmetries:
group-equivariant networks (\Cref{ch:groups}) generalize equivariance from
translations to finite and continuous groups such as $D_4$, $SO(2)$, and
$SO(3)$, addressing the symmetry limitations; graph neural networks
(\Cref{ch:graphs}) overcome the topological limitations through generalization
to irregular relational domains; convolutions on manifolds
(\Cref{ch:geodesics}) extend to curved surfaces; and gauge convolutions
(\Cref{ch:gauges}) introduce equivariance to local gauge transformations on
fibre bundles, the most abstract generalization. This systematic progression
from restricted equivariance towards full generality is the central narrative of
Part II and exemplifies the principles of geometric deep learning in action.

\section{Harmonic analysis and spectral convolutions}
\label{sec:grids:harmonic}

While the geometric extensions of CNNs address the limitations of symmetry,
topology, and locality through the direct generalization of convolutional
operations, harmonic analysis provides a dual and complementary perspective that
reveals the deep mathematical structure underlying all these architectures. This
section develops the theoretical framework connecting classical convolutions
with their spectral generalizations, from the foundations of Fourier analysis to
modern applications on irregular domains.

The spatial--spectral duality is one of the most fundamental principles in
applied mathematics: while convolution operations are naturally carried out in
the spatial domain, their algebraic structure is best understood in the
frequency domain. The convolution theorem,
$\Ft[f \conv g] = \Ft[f] \cdot \Ft[g]$, reveals that the complex convolution
operation in space reduces to a simple pointwise multiplication in frequency.
This perspective not only provides computational advantages through algorithms
such as the Fast Fourier Transform (FFT), but also offers deep insights into
what features neural networks learn and how more efficient architectures may be
designed.

In the context of geometric deep learning, harmonic analysis acquires
fundamental importance by providing the mathematical tools needed to generalize
convolutions to non-Euclidean domains. The representation theory of groups,
developed within the framework of harmonic analysis (\Cref{ch:foundations}),
allows the design of convolutional layers that automatically preserve the
specific symmetries of the data domain. Likewise, the spectral analysis of
graphs through the Laplacian provides the theoretical basis for spectral graph
networks (\Cref{ch:graphs}), while wavelets offer natural hierarchical
representations aligned with modern deep architectures.

\subsection{Foundations of Fourier analysis}
\label{sec:grids:fourier-foundations}

Fourier analysis is one of the most powerful mathematical tools for the study of
periodic phenomena and frequency analysis. Its application in the context of
convolutional neural networks enables a deep understanding of the spectral
properties of filters and convolutional operations.

\begin{definition}[Continuous Fourier transform]
\label{def:grids:fourier}
Let $f \colon \R^d \to \C$ be an integrable function. The \emph{Fourier
transform} of $f$ is defined as
\[
  \Ft[f](\omega) = \hat{f}(\omega)
   = \int_{\R^d} f(x)\, e^{-2\pi i\, \omega \cdot x}\, \mathrm{d}x,
\]
where $\omega \in \R^d$ is the frequency variable and
$\omega \cdot x = \sum_{j=1}^d \omega_j x_j$ is the Euclidean inner product.
\end{definition}

\paragraph{Existence conditions.}
The Fourier transform is well defined for functions in $L^1(\R^d)$ (absolutely
integrable functions). For functions in $L^2(\R^d)$, the transform is defined as
a mean-square limit.

\begin{definition}[Inverse Fourier transform]
\label{def:grids:fourier-inverse}
If $\hat{f} \in L^1(\R^d)$, the \emph{inverse Fourier transform} is defined as
\[
  \Ft^{-1}[\hat{f}](x) = f(x)
   = \int_{\R^d} \hat{f}(\omega)\, e^{2\pi i\, \omega \cdot x}\, \mathrm{d}\omega.
\]
\end{definition}

\begin{theorem}[Fourier inversion theorem]
\label{thm:grids:inversion}
If $f, \hat{f} \in L^1(\R^d)$, then
\[
  f(x) = \Ft^{-1}[\Ft[f]](x)
   = \int_{\R^d} \hat{f}(\omega)\, e^{2\pi i\, \omega \cdot x}\, \mathrm{d}\omega
\]
for almost every $x \in \R^d$. For the proof, see \citep[Theorem~8.26]{Folland1999}.
\end{theorem}

The complex exponentials $e^{2\pi i\,\omega\cdot x}$ are exactly the
\emph{eigenfunctions} of translation: shifting the argument multiplies the
exponential by a phase. This is what makes Fourier analysis the natural language
of translation symmetry, and it foreshadows the central role of the Laplacian
eigenbasis on graphs (\Cref{ch:graphs}), which plays the same part on irregular
domains.

\begin{definition}[Convolution on $\R^d$]
\label{def:grids:convolucion-continua}
For functions $f, g \colon \R^d \to \C$, the \emph{convolution} is defined as
\[
  (f \conv g)(x) = \int_{\R^d} f(y)\, g(x - y)\, \mathrm{d}y
   = \int_{\R^d} f(x - y)\, g(y)\, \mathrm{d}y,
\]
when the integral converges.
\end{definition}

\paragraph{Concrete examples of Fourier transforms.}
The following examples illustrate the fundamental properties of the Fourier
transform in contexts relevant to the analysis of CNNs.

\begin{example}[Transform of the Gaussian function]
\label{ex:grids:fourier-gaussian}
For the Gaussian function $f(x) = e^{-\alpha x^2}$ with $\alpha > 0$,
\[
  \Ft[e^{-\alpha x^2}](\omega) = \sqrt{\frac{\pi}{\alpha}}\; e^{-\pi^2 \omega^2/\alpha}.
\]
Notable properties: the transform of a Gaussian is another Gaussian; there is a
dual localization (a narrow function corresponds to a wide transform); it
minimizes the time--frequency uncertainty principle; and it is fundamental in
smoothing filters for CNNs.
\end{example}

\begin{example}[Transform of the step function]
\label{ex:grids:fourier-step}
For the unit step function $\mathbf{1}_{[0,1]}(x)$,
\[
  \Ft[\mathbf{1}_{[0,1]}](\omega)
   = \frac{e^{-2\pi i \omega} - 1}{-2\pi i \omega}
   = \frac{\sin(\pi\omega)}{\pi\omega}\, e^{-\pi i \omega}.
\]
This represents averaging filters (constant kernels); the sinc function in
frequency indicates low-pass behaviour, and the oscillations in frequency are
due to discontinuities in the spatial domain.
\end{example}

\begin{example}[Transform of the sinc function]
\label{ex:grids:fourier-sinc}
For the sinc function $f(x) = \dfrac{\sin(\pi x)}{\pi x}$,
\[
  \Ft[\sinc(x)](\omega) = \mathbf{1}_{[-1/2,\,1/2]}(\omega).
\]
This step $\leftrightarrow$ sinc duality is fundamental: an ideal low-pass
filter in frequency corresponds to a sinc response in space; it is impossible to
implement exactly (infinite support); and it inspires the design of approximate
filters in spectral CNNs.
\end{example}

\begin{example}[Transform of the Dirac delta]
\label{ex:grids:fourier-delta}
For the delta distribution $\delta(x)$,
\[
  \Ft[\delta](\omega) = 1 \quad \text{and} \quad \Ft[1](\omega) = \delta(\omega).
\]
A spatial delta has a flat spectrum (all frequencies present); a spatial
constant is a frequency impulse (zero frequency only). This embodies the
localization principle in the spectral analysis of filters.
\end{example}

\begin{figure}[t]
\centering
\begin{tikzpicture}[scale=0.85]
\node[font=\bfseries, text=gdlblue!70!black] at (3.2, 1.0) {Fourier transform pairs};
% Row 1: Gaussian <-> Gaussian
\begin{scope}[shift={(0, -1.2)}]
  \draw[-{Stealth[length=4pt]}, gdlgray!60] (-1.7, 0) -- (1.9, 0)
    node[right, font=\tiny, text=gdlgray!70!black] {$t$};
  \draw[-{Stealth[length=4pt]}, gdlgray!60] (0, -0.2) -- (0, 1.6);
  \draw[gdlblue!70, thick] plot[smooth, domain=-1.6:1.6, samples=40] (\x, {1.3*exp(-\x*\x)});
  \node[font=\tiny, text=gdlblue!70!black] at (1.1, 1.2) {$e^{-t^2}$};
  \node[font=\tiny\itshape, text=gdlgray!70!black, anchor=north] at (0, -0.4) {Gaussian};
\end{scope}
\draw[-{Stealth[length=5pt]}, gdlpurple!70, double, double distance=1pt]
  (2.4, -0.7) -- (3.6, -0.7) node[midway, above, font=\tiny, text=gdlpurple!70!black] {$\Ft$};
\begin{scope}[shift={(6, -1.2)}]
  \draw[-{Stealth[length=4pt]}, gdlgray!60] (-1.7, 0) -- (1.9, 0)
    node[right, font=\tiny, text=gdlgray!70!black] {$\omega$};
  \draw[-{Stealth[length=4pt]}, gdlgray!60] (0, -0.2) -- (0, 1.6);
  \draw[gdlblue!70, thick] plot[smooth, domain=-1.6:1.6, samples=40] (\x, {1.3*exp(-\x*\x)});
  \node[font=\tiny, text=gdlblue!70!black] at (1.1, 1.2) {$e^{-\omega^2}$};
  \node[font=\tiny\itshape, text=gdlgray!70!black, anchor=north] at (0, -0.4) {Gaussian};
\end{scope}
% Row 2: Rect <-> Sinc
\begin{scope}[shift={(0, -3.4)}]
  \draw[-{Stealth[length=4pt]}, gdlgray!60] (-1.7, 0) -- (1.9, 0)
    node[right, font=\tiny, text=gdlgray!70!black] {$t$};
  \draw[-{Stealth[length=4pt]}, gdlgray!60] (0, -0.2) -- (0, 1.6);
  \draw[gdlorange!70, thick] (-1.6,0) -- (-0.7,0) -- (-0.7,1.1) -- (0.7,1.1) -- (0.7,0) -- (1.6,0);
  \node[font=\tiny, text=gdlorange!70!black] at (0, 1.4) {$\operatorname{rect}(t)$};
  \node[font=\tiny\itshape, text=gdlgray!70!black, anchor=north] at (0, -0.4) {Rectangular};
\end{scope}
\draw[-{Stealth[length=5pt]}, gdlpurple!70, double, double distance=1pt]
  (2.4, -2.9) -- (3.6, -2.9) node[midway, above, font=\tiny, text=gdlpurple!70!black] {$\Ft$};
\begin{scope}[shift={(6, -3.4)}]
  \draw[-{Stealth[length=4pt]}, gdlgray!60] (-1.7, 0) -- (1.9, 0)
    node[right, font=\tiny, text=gdlgray!70!black] {$\omega$};
  \draw[-{Stealth[length=4pt]}, gdlgray!60] (0, -0.4) -- (0, 1.6);
  \draw[gdlorange!70, thick] plot[smooth, domain=-1.6:-0.05, samples=40] (\x, {1.1*sin(180*\x)/(\x*3.14159)});
  \draw[gdlorange!70, thick] plot[smooth, domain=0.05:1.6, samples=40] (\x, {1.1*sin(180*\x)/(\x*3.14159)});
  \fill[gdlorange!70] (0,1.1) circle (1pt);
  \node[font=\tiny, text=gdlorange!70!black] at (1.15, 1.2) {$\operatorname{sinc}(\omega)$};
  \node[font=\tiny\itshape, text=gdlgray!70!black, anchor=north] at (0, -0.55) {Sinc};
\end{scope}
% Row 3: Delta <-> Constant
\begin{scope}[shift={(0, -5.6)}]
  \draw[-{Stealth[length=4pt]}, gdlgray!60] (-1.7, 0) -- (1.9, 0)
    node[right, font=\tiny, text=gdlgray!70!black] {$t$};
  \draw[-{Stealth[length=4pt]}, gdlgray!60] (0, -0.2) -- (0, 1.6);
  \draw[gdlred!70, very thick, -{Stealth[length=4pt]}] (0,0) -- (0,1.3);
  \node[font=\tiny, text=gdlred!70!black, anchor=west] at (0.1, 1.3) {$\delta(t)$};
  \node[font=\tiny\itshape, text=gdlgray!70!black, anchor=north] at (0, -0.4) {Delta};
\end{scope}
\draw[-{Stealth[length=5pt]}, gdlpurple!70, double, double distance=1pt]
  (2.4, -5.1) -- (3.6, -5.1) node[midway, above, font=\tiny, text=gdlpurple!70!black] {$\Ft$};
\begin{scope}[shift={(6, -5.6)}]
  \draw[-{Stealth[length=4pt]}, gdlgray!60] (-1.7, 0) -- (1.9, 0)
    node[right, font=\tiny, text=gdlgray!70!black] {$\omega$};
  \draw[-{Stealth[length=4pt]}, gdlgray!60] (0, -0.2) -- (0, 1.6);
  \draw[gdlred!70, thick] (-1.6, 0.9) -- (1.6, 0.9);
  \node[font=\tiny, text=gdlred!70!black, anchor=south west] at (1.0, 1.0) {$1$};
  \node[font=\tiny\itshape, text=gdlgray!70!black, anchor=north] at (0, -0.4) {Constant};
\end{scope}
\node[font=\scriptsize, text=gdlpurple!70!black, anchor=north] at (3.2, -7.0)
  {Localization in $t$ $\longleftrightarrow$ delocalization in $\omega$
   \quad (uncertainty principle)};
\end{tikzpicture}
\caption[Fourier transform pairs]{Fourier transform pairs: correspondence
between signals in the time domain ($t$) and their representations in the
frequency domain ($\omega$). The localization--delocalization duality that
underlies the uncertainty principle is apparent.}
\label{fig:grids:fourier-pairs}
\end{figure}

The space--frequency duality of these pairs has a direct architectural meaning.
Gaussians are the basis for Gabor wavelets (time--frequency analysis); step
functions are related to Haar wavelets (edge features); and sinc functions
underlie Daubechies wavelets (perfect reconstruction). The ability to analyse
multiple scales simultaneously makes wavelets especially attractive for
hierarchical CNN architectures, as developed in
\Cref{sec:grids:wavelets}.

\subsection{Fundamental properties of harmonic analysis}
\label{sec:grids:harmonic-properties}

\begin{theorem}[Convolution theorem in the Fourier domain]
\label{thm:grids:conv-theorem}
For functions $f, g \colon \R^d \to \C$,
\[
  \Ft[f \conv g] = \Ft[f] \cdot \Ft[g],
\]
where $\Ft$ denotes the Fourier transform.
\end{theorem}

\begin{proof}
Let $f, g \in L^1(\R^d)$. Applying Fubini's theorem,
\begin{align}
\Ft[f \conv g](\omega)
  &= \int_{\R^d} (f \conv g)(x)\, e^{-2\pi i\omega \cdot x}\, \mathrm{d}x \nonumber \\
  &= \int_{\R^d} \Big[\int_{\R^d} f(y)\,g(x-y)\, \mathrm{d}y\Big]
       e^{-2\pi i\omega \cdot x}\, \mathrm{d}x \nonumber \\
  &= \int_{\R^d} f(y) \Big[\int_{\R^d} g(x-y)\, e^{-2\pi i\omega \cdot x}\, \mathrm{d}x\Big] \mathrm{d}y \nonumber \\
  &= \int_{\R^d} f(y)\, e^{-2\pi i\omega \cdot y}\, \Ft[g](\omega)\, \mathrm{d}y
   = \Ft[f](\omega) \cdot \Ft[g](\omega),
\end{align}
where in the third step we used the change of variable $z = x - y$ and the
definition of $\Ft[g]$.
\end{proof}

The theorem turns the global sliding operation of convolution into an
independent rescaling of each frequency component. It gives a clean
interpretation of what a filter does: a kernel acts as a frequency mask,
amplifying some frequencies and attenuating others. It exposes the meaning of
the hand-designed kernels of \Cref{ex:grids:sobel}: a smoothing (Gaussian)
kernel is a low-pass filter, while a Sobel or Laplacian edge detector is a
high-pass filter. Learning a kernel is thus learning a frequency response.

\begin{corollary}[Computational efficiency via the FFT]
\label{cor:grids:fft}
The convolution can be computed in $O(n \log n)$ using
\[
  f \conv g = \Ft^{-1}[\Ft[f] \cdot \Ft[g]].
\]
\end{corollary}

\begin{proof}
By the convolution theorem, $\Ft[f \conv g] = \Ft[f] \cdot \Ft[g]$, hence
$f \conv g = \Ft^{-1}[\Ft[f] \cdot \Ft[g]]$. For discrete signals of length $n$,
the FFT computes $\Ft[f]$ and $\Ft[g]$ each in $O(n \log n)$ operations, the
pointwise multiplication $\Ft[f] \cdot \Ft[g]$ requires $O(n)$ operations, and
the inverse transform $\Ft^{-1}$ has the same cost $O(n \log n)$. The total cost
is $O(n \log n)$, compared with the $O(n^2)$ of the direct computation.
\end{proof}

\begin{theorem}[Plancherel identity]
\label{thm:grids:plancherel}
For functions $f, g \in L^2(\R^d)$, the Fourier transform preserves the inner
product:
\[
  \inner{f}{g}_{L^2} = \inner{\hat{f}}{\hat{g}}_{L^2}.
\]
In particular, $\norm{\hat{f}}_{L^2} = \norm{f}_{L^2}$ (isometry). See
\citep[Theorem~8.29]{Folland1999}.
\end{theorem}

The Plancherel identity guarantees that operations in the frequency domain
preserve the energy of signals, which is fundamental for the numerical stability
of spectral algorithms in CNNs.

\begin{definition}[Discrete Fourier transform]
\label{def:grids:dft}
For a finite sequence $\{x_n\}_{n=0}^{N-1}$, the \emph{discrete Fourier
transform} (DFT) is defined as
\[
  X_k = \sum_{n=0}^{N-1} x_n\, e^{-2\pi i kn/N}, \qquad k = 0, 1, \ldots, N-1.
\]
\end{definition}

\begin{definition}[Fast Fourier transform]
\label{def:grids:fft}
The \emph{fast Fourier transform} (FFT) algorithm computes the DFT in time
$O(N \log N)$ instead of $O(N^2)$ through recursive factorization:
\[
  X_k = \sum_{n=0}^{N/2-1} x_{2n}\, e^{-2\pi i kn/(N/2)}
   + e^{-2\pi i k/N} \sum_{n=0}^{N/2-1} x_{2n+1}\, e^{-2\pi i kn/(N/2)}.
\]
\end{definition}

\begin{theorem}[Nyquist--Shannon sampling theorem]
\label{thm:grids:nyquist}
Let $f \colon \R \to \C$ be a band-limited function with $\hat{f}(\omega) = 0$
for $\abs{\omega} > W$. Then $f$ is completely determined by its samples at
intervals of $1/(2W)$:
\[
  f(t) = \sum_{n=-\infty}^{\infty} f\!\Big(\frac{n}{2W}\Big)
    \sinc\!\Big(2W\Big(t - \frac{n}{2W}\Big)\Big),
\]
where $\sinc(x) = \sin(\pi x)/(\pi x)$. This fundamental result was established
by \citet{Shannon1948}.
\end{theorem}

\paragraph{Nyquist frequency and aliasing.}
The Nyquist frequency $f_N = W$ is the maximum frequency that can be represented
without aliasing. In CNNs, subsampling (pooling) can introduce aliasing if
appropriate anti-aliasing filters are not applied.

\begin{definition}[Fourier series]
\label{def:grids:fourier-series}
For a periodic function $f \colon \R \to \C$ with period $T$, the Fourier series
is
\[
  f(t) = \sum_{n=-\infty}^{\infty} c_n\, e^{2\pi i n t/T},
\]
where the coefficients are $c_n = \dfrac{1}{T} \int_0^T f(t)\, e^{-2\pi i n t/T}\, \mathrm{d}t$.
\end{definition}

\begin{definition}[Short-time Fourier transform]
\label{def:grids:stft}
To analyse non-stationary signals, the \emph{short-time Fourier transform}
(STFT) is defined as
\[
  \mathrm{STFT}_f(t, \omega)
   = \int_{-\infty}^{\infty} f(\tau)\, w(\tau - t)\, e^{-2\pi i \omega \tau}\, \mathrm{d}\tau,
\]
where $w$ is a window function (typically Gaussian or Hann).
\end{definition}

\paragraph{Time--frequency uncertainty principle.}
The STFT is subject to the uncertainty principle,
$\Delta t \cdot \Delta \omega \geq \frac{1}{4\pi}$, which implies a fundamental
trade-off between temporal and frequency resolution.

\begin{example}[Fast convolution via FFT: complexity analysis]
\label{ex:grids:fft-complexity}
Consider the convolution of an image $f \in \R^{N \times N}$ with a kernel
$g \in \R^{K \times K}$. The direct method,
\[
  (f \conv g)[i,j] = \sum_{u=0}^{K-1} \sum_{v=0}^{K-1} f[i-u,\, j-v] \cdot g[u,v],
\]
has complexity $O(N^2 K^2)$. The spectral method computes
$\hat{f} = \mathrm{FFT2D}(f)$ and $\hat{g} = \mathrm{FFT2D}(\mathrm{pad}(g,N))$,
each in $O(N^2 \log N)$, then $\hat{h} = \hat{f} \odot \hat{g}$ in $O(N^2)$, and
finally $h = \mathrm{IFFT2D}(\hat{h})$ in $O(N^2 \log N)$, for a total of
$O(N^2 \log N)$. The spectral method is more efficient when $K > O(\log N)$.
\end{example}

\subsection{Harmonic analysis on groups}
\label{sec:grids:harmonic-groups}

Harmonic analysis on groups generalizes the classical Fourier transform to more
general spaces with group structure, providing the mathematical tools
fundamental to geometric convolutions in neural networks. The full theory of
group representations --- the definition of a representation, irreducible
representations, and the Peter--Weyl theorem --- is developed in
\Cref{ch:foundations}. Here we briefly recall the specific harmonic-analysis
tools we use for CNNs.

We recall that the \emph{Haar measure} of a locally compact group $G$ is the
unique Borel measure (up to a constant) invariant under left translations:
$\mu(gE) = \mu(E)$ for all $g \in G$. Haar's theorem guarantees that the Haar
measure exists and is unique up to a multiplicative constant. For finite groups,
the Haar measure is the counting measure.

The \emph{character} of a representation $\rho \colon G \to GL(\C^d)$ is the
function $\chi(g) = \tr(\rho(g))$, and the irreducible characters satisfy the
orthogonality relation $\frac{1}{\abs{G}} \sum_{g \in G} \chi_i(g)
\overline{\chi_j(g)} = \delta_{ij}$. From these tools the Fourier transform on
finite groups is defined: for $f \colon G \to \C$ and an irreducible
representation $\rho_i$ of dimension $d_i$,
\begin{equation}
  \hat{f}(\rho_i) = \sum_{g\in G} f(g)\,\rho_i(g) \in \C^{d_i \times d_i},
  \label{eq:grids:fourier-finite-group}
\end{equation}
with inversion $f(g) = \frac{1}{\abs{G}} \sum_i d_i\, \tr\big(\hat{f}(\rho_i)\,
\rho_i(g^{-1})\big)$. The fundamental result for CNNs is the \textbf{convolution
theorem}: the Fourier transform converts group convolution into a matrix
product,
\begin{equation}
  \widehat{f \conv_G h}(\rho_i) = \hat{f}(\rho_i)\, \hat{h}(\rho_i),
  \label{eq:grids:conv-fourier-groups}
\end{equation}
generalizing the classical DFT result.

\begin{example}[Cyclic group $\Z_n$]
\label{ex:grids:cyclic-group}
For the cyclic group of order $n$, the irreducible representations are
one-dimensional:
\[
  \rho_k(j) = e^{-2\pi i jk/n}, \qquad k = 0, 1, \ldots, n-1.
\]
The Fourier transform reduces to the classical DFT,
$\hat{f}(k) = \sum_{j=0}^{n-1} f(j)\, e^{-2\pi i jk/n}$.
\end{example}

\begin{example}[Dihedral group $D_n$]
\label{ex:grids:dihedral-group}
For the dihedral group of order $2n$ (symmetries of an $n$-gon), the
representations include trivial representations ($\rho(g) = 1$ for all $g$), sign
representations ($\rho(r^j) = 1$, $\rho(sr^j) = -1$), and two-dimensional
representations for the combined rotations and reflections.
\end{example}

\begin{definition}[Fourier transform on compact groups]
\label{def:grids:fourier-compact}
For a compact group $G$ and a function $f \in L^2(G)$, the Fourier transform is
defined through the Peter--Weyl theorem:
\[
  f = \sum_{\lambda \in \widehat{G}} d_\lambda\, \tr\big(\hat{f}(\lambda)\, U^\lambda\big),
\]
where $\widehat{G}$ indexes the irreducible representations.
\end{definition}

\begin{definition}[Circular harmonics on $SO(2)$]
\label{def:grids:circular-harmonics}
For the group of planar rotations $SO(2) \cong S^1$, the irreducible
representations are $U^n(\theta) = e^{in\theta}$, $n \in \Z$. A function
$f \colon S^1 \to \C$ expands as
$f(\theta) = \sum_{n=-\infty}^{\infty} c_n\, e^{in\theta}$.
\end{definition}

\begin{definition}[Spherical harmonics on $SO(3)$]
\label{def:grids:spherical-harmonics}
For the group of three-dimensional rotations $SO(3)$, functions on the sphere
$S^2$ expand in spherical harmonics:
\[
  f(\theta, \phi) = \sum_{l=0}^{\infty} \sum_{m=-l}^{l} c_{lm}\, Y_l^m(\theta, \phi),
\]
where $Y_l^m$ are the spherical harmonics of degree $l$ and order $m$.
\end{definition}

The Peter--Weyl theorem (\Cref{ch:foundations}) provides the theoretical basis
for the decomposition into channels in group-equivariant CNNs, the design of
equivariant filters on spheres, and the efficient parametrization of
convolutions on Lie groups --- topics developed in detail in \Cref{ch:groups}.

\subsection{Spectral analysis on graphs}
\label{sec:grids:spectral-graphs}

Spectral analysis on graphs extends the ideas of Fourier analysis to irregular
domains, providing the mathematical basis for spectral convolutions on graph
neural networks. The theoretical foundations of the graph Laplacian, including
its definition and spectral properties, are developed in \Cref{ch:graphs}; here
we summarize only the consequences needed to relate grid convolution to its
spectral generalization. For the Laplacian matrix $L = D - A$ and its normalized
variant $\mathcal{L} = D^{-1/2} L\, D^{-1/2}$, we assume familiarity with the
spectral decomposition.

\begin{theorem}[Spectral theorem for the Laplacian]
\label{thm:grids:spectral-laplacian}
Let $L$ be the Laplacian matrix of a connected graph. Then:
\begin{enumerate}
  \item The eigenvalues satisfy $0 = \lambda_0 < \lambda_1 \leq \lambda_2 \leq
    \cdots \leq \lambda_{n-1}$.
  \item The eigenvector associated with $\lambda_0 = 0$ is
    $\phi_0 = \frac{1}{\sqrt{n}} \mathbf{1}$.
  \item The eigenvectors $\{\phi_k\}_{k=0}^{n-1}$ form an orthonormal basis of
    $\R^n$.
  \item The multiplicity of $\lambda_0$ equals the number of connected
    components.
\end{enumerate}
For the proof, see \citep[Chapters~1--2]{Chung1997} and \Cref{ch:graphs}.
\end{theorem}

The graph Fourier transform (\Cref{ch:graphs}) projects signals onto the
eigenvectors of the Laplacian: $\hat{f}(\lambda_k) = \inner{\mathbf{f}}{\phi_k}$,
with matrix form $\hat{f} = U^T f$. From this transform the spectral
convolutions are defined.

\begin{definition}[Spectral convolution on graphs]
\label{def:grids:spectral-conv-graphs}
For a function $f \colon V \to \R$ and a spectral filter
$g \colon \R_+ \to \R$,
\[
  (g \star f)(v_i) = \sum_{k=0}^{n-1} g(\lambda_k)\, \hat{f}(\lambda_k)\, \phi_k(v_i).
\]
In matrix form, $g \star f = \Phi\, g(\Lambda)\, \Phi^T f$, where
$\Phi = [\phi_0, \ldots, \phi_{n-1}]$ and
$\Lambda = \diag(\lambda_0, \ldots, \lambda_{n-1})$.
\end{definition}

\begin{theorem}[Frequency interpretation]
\label{thm:grids:freq-interp-graphs}
The eigenvalues of the Laplacian admit a frequency interpretation: low
frequencies ($\lambda_k$ small) correspond to eigenvectors that vary smoothly
over the graph, high frequencies ($\lambda_k$ large) correspond to eigenvectors
that oscillate rapidly between adjacent vertices, and low-pass filters satisfy
$g(\lambda) \approx 1$ for small $\lambda$ and $g(\lambda) \approx 0$ for large
$\lambda$.
\end{theorem}

\begin{proof}
The interpretation is justified through the Rayleigh quotient. For a signal
$\mathbf{f} \in \R^n$ over the graph, the total variation is
$\mathbf{f}^T L \mathbf{f} = \sum_{(i,j) \in E} w_{ij}(f_i - f_j)^2$. The
$k$-th eigenvector $\phi_k$ minimizes this total variation subject to
orthogonality with the previous eigenvectors:
$\lambda_k = \min_{\mathbf{f} \perp \phi_0, \ldots, \phi_{k-1}}
\frac{\mathbf{f}^T L \mathbf{f}}{\mathbf{f}^T \mathbf{f}}$. Thus eigenvectors
with small eigenvalues minimize the variation between adjacent nodes (smooth,
low-frequency variation), while large eigenvalues correspond to eigenvectors
with high variation (rapid oscillations, high frequency).
\end{proof}

To make spectral convolutions computationally efficient, filters are
approximated by polynomials, avoiding the explicit diagonalization of the
Laplacian. Using Chebyshev polynomials $T_k$ on the interval $[-1, 1]$, a filter
is written as $g(\tilde{\lambda}) = \sum_{k=0}^K \theta_k\, T_k(\tilde{\lambda})$
with $\tilde{\lambda} = 2\lambda/\lambda_{\max} - 1$, and the action $g(L)f$ is
computed in $O(K\abs{E})$ through the recurrence
$\bar{x}_k = 2\tilde{L}\bar{x}_{k-1} - \bar{x}_{k-2}$, without diagonalizing $L$.
These spectral methods on graphs are developed fully in \Cref{ch:graphs}; here
they serve to show that grid convolution is the regular special case of a
construction that survives the loss of the grid.

\subsection{Wavelet bases and multiresolution analysis}
\label{sec:grids:wavelets}

Wavelets provide an alternative to the Fourier transform that permits
simultaneous time--frequency analysis, overcoming the limitations of the
uncertainty principle. Their application in neural networks offers natural
hierarchical representations.

\begin{definition}[Mother wavelet]
\label{def:grids:mother-wavelet}
A function $\psi \in L^2(\R)$ is a \emph{mother wavelet} if it satisfies the
admissibility condition
\[
  C_\psi = \int_{-\infty}^{\infty} \frac{\abs{\hat{\psi}(\omega)}^2}{\abs{\omega}}\, \mathrm{d}\omega < \infty,
\]
where $\hat{\psi}$ is the Fourier transform of $\psi$.
\end{definition}

The condition $C_\psi < \infty$ implies that $\hat{\psi}(0) = 0$, which means
$\int_{-\infty}^{\infty} \psi(t)\, \mathrm{d}t = 0$. This guarantees that
wavelets have zero mean and provide local analysis.

\begin{definition}[Continuous wavelet transform]
\label{def:grids:cwt}
For a function $f \in L^2(\R)$ and a mother wavelet $\psi$, the \emph{continuous
wavelet transform} is defined as
\[
  W_f(a, b) = \frac{1}{\sqrt{\abs{a}}} \int_{-\infty}^{\infty} f(t)\,
    \psi^*\!\Big(\frac{t-b}{a}\Big)\, \mathrm{d}t,
\]
where $a \in \R\setminus\{0\}$ is the scale parameter and $b \in \R$ is the
translation parameter.
\end{definition}

\begin{theorem}[Wavelet reconstruction formula]
\label{thm:grids:wavelet-reconstruction}
If $\psi$ satisfies the admissibility condition, then
\[
  f(t) = \frac{1}{C_\psi} \int_{-\infty}^{\infty} \int_{-\infty}^{\infty}
    W_f(a, b)\, \psi_{a,b}(t)\, \frac{\mathrm{d}a\, \mathrm{d}b}{a^2},
\]
where $\psi_{a,b}(t) = \frac{1}{\sqrt{\abs{a}}}\, \psi\big(\frac{t-b}{a}\big)$.
See \Cref{ch:appendix} for the complete proof.
\end{theorem}

\begin{definition}[Multiresolution analysis]
\label{def:grids:mra}
A \emph{multiresolution analysis} (MRA) of $L^2(\R)$ is a sequence of closed
subspaces $\{V_j\}_{j \in \Z}$ satisfying:
\begin{enumerate}
  \item \textbf{Inclusion:} $\cdots \subset V_{-1} \subset V_0 \subset V_1 \subset \cdots$.
  \item \textbf{Density:} $\overline{\bigcup_{j \in \Z} V_j} = L^2(\R)$.
  \item \textbf{Separation:} $\bigcap_{j \in \Z} V_j = \{0\}$.
  \item \textbf{Scaling:} $f(x) \in V_j \Leftrightarrow f(2x) \in V_{j+1}$.
  \item \textbf{Translation:} $f(x) \in V_0 \Rightarrow f(x-k) \in V_0$ for $k \in \Z$.
  \item \textbf{Riesz basis:} there exists $\phi \in V_0$ such that
    $\{\phi(x-k)\}_{k \in \Z}$ is a Riesz basis of $V_0$.
\end{enumerate}
\end{definition}

\begin{definition}[Scaling function]
\label{def:grids:scaling-function}
The function $\phi$ in condition~6 of the MRA is called the \emph{scaling
function} or \emph{father wavelet}. It satisfies the refinement equation
\[
  \phi(x) = \sqrt{2} \sum_{k \in \Z} h_k\, \phi(2x - k),
\]
where $\{h_k\}$ are the scaling coefficients.
\end{definition}

\begin{definition}[Orthogonal wavelet]
\label{def:grids:orthogonal-wavelet}
Given an MRA with scaling function $\phi$, the associated \emph{orthogonal
wavelet} is defined as
\[
  \psi(x) = \sqrt{2} \sum_{k \in \Z} g_k\, \phi(2x - k),
\]
where $g_k = (-1)^k h_{1-k}$ are related to the scaling coefficients.
\end{definition}

\begin{theorem}[Wavelet decomposition]
\label{thm:grids:wavelet-decomposition}
Every function $f \in L^2(\R)$ can be decomposed as
\[
  f(x) = \sum_{k \in \Z} c_{J,k}\, \phi_{J,k}(x)
   + \sum_{j=J}^{\infty} \sum_{k \in \Z} d_{j,k}\, \psi_{j,k}(x),
\]
where $\phi_{j,k}(x) = 2^{j/2} \phi(2^j x - k)$ and
$\psi_{j,k}(x) = 2^{j/2} \psi(2^j x - k)$. See \Cref{ch:appendix} for the proof.
\end{theorem}

\begin{figure}[t]
\centering
\begin{tikzpicture}[
  >=Stealth, scale=0.9, transform shape,
  filter/.style={draw=#1!70, fill=#1!15, minimum width=0.9cm, minimum height=0.65cm,
    font=\small, rounded corners=2pt, inner sep=2pt},
  downsample/.style={draw=gdlgray!70, fill=gdlgray!15, minimum width=0.65cm,
    minimum height=0.65cm, font=\small, rounded corners=2pt, inner sep=2pt},
  coeff/.style={draw=#1!60, fill=#1!15, minimum width=1.5cm, minimum height=0.65cm,
    font=\small, rounded corners=3pt, inner sep=3pt},
  arrowline/.style={->, gdlgray!80, thick},
]
\def\rowsep{2.4}
\def\collow{-1.6}
\def\colhigh{3.8}
\node[draw=gdlgray!60, fill=gdlgray!10, minimum width=2.6cm, minimum height=1.0cm,
      rounded corners=3pt] (fbox) at (1.1, 0) {};
\begin{scope}
  \clip ([xshift=3pt, yshift=3pt]fbox.south west) rectangle ([xshift=-3pt, yshift=-3pt]fbox.north east);
  \draw[gdlred!70, line width=1.0pt]
    plot[domain=-1.15:1.15, samples=100, smooth]
    ({1.1 + \x}, {0.3*sin(4*\x r) + 0.13*sin(11*\x r) + 0.09*cos(18*\x r)});
\end{scope}
\node[above=2pt of fbox, font=\bfseries\small] {Signal $f \in V_J$};
\def\lva{-2.3}
\draw[arrowline] (fbox.south) -- ++(0, -0.3) coordinate (splitA);
\draw[gdlgray!80, thick] (splitA) -| (\collow, {\lva + 0.55});
\draw[arrowline] (\collow, {\lva + 0.55}) -- (\collow, {\lva + 0.33});
\draw[gdlgray!80, thick] (splitA) -| (\colhigh, {\lva + 0.55});
\draw[arrowline] (\colhigh, {\lva + 0.55}) -- (\colhigh, {\lva + 0.33});
\node[filter=gdlblue] (hA) at (\collow, \lva) {$h$};
\node[above=1pt of hA, font=\scriptsize, text=gdlblue!70!black] {low-pass};
\node[downsample, right=0.3cm of hA] (dsA) {$\downarrow\!2$};
\draw[arrowline] (hA) -- (dsA);
\node[coeff=gdlblue, right=0.45cm of dsA] (cJ1) {$c_{J-1,k}$};
\draw[arrowline] (dsA) -- (cJ1);
\node[below=2pt of cJ1, font=\scriptsize, text=gdlblue!70!black] {$\in V_{J-1}$};
\node[filter=gdlorange] (gA) at (\colhigh, \lva) {$g$};
\node[above=1pt of gA, font=\scriptsize, text=gdlorange!70!black] {high-pass};
\node[downsample, right=0.3cm of gA] (dsAgA) {$\downarrow\!2$};
\draw[arrowline] (gA) -- (dsAgA);
\node[coeff=gdlorange, right=0.45cm of dsAgA] (dJ1) {$d_{J-1,k}$};
\draw[arrowline] (dsAgA) -- (dJ1);
\node[below=2pt of dJ1, font=\scriptsize, text=gdlorange!70!black] {$\in W_{J-1}$};
\pgfmathsetmacro{\lvb}{\lva - \rowsep}
\draw[arrowline] (cJ1.south) -- ++(0, -0.5) coordinate (splitB);
\draw[gdlgray!80, thick] (splitB) -| (\collow, {\lvb + 0.55});
\draw[arrowline] (\collow, {\lvb + 0.55}) -- (\collow, {\lvb + 0.33});
\draw[gdlgray!80, thick] (splitB) -| (\colhigh, {\lvb + 0.55});
\draw[arrowline] (\colhigh, {\lvb + 0.55}) -- (\colhigh, {\lvb + 0.33});
\node[filter=gdlblue] (hB) at (\collow, \lvb) {$h$};
\node[downsample, right=0.3cm of hB] (dsB) {$\downarrow\!2$};
\draw[arrowline] (hB) -- (dsB);
\node[coeff=gdlblue, right=0.45cm of dsB] (cJ2) {$c_{J-2,k}$};
\draw[arrowline] (dsB) -- (cJ2);
\node[below=2pt of cJ2, font=\scriptsize, text=gdlblue!70!black] {$\in V_{J-2}$};
\node[filter=gdlorange] (gB) at (\colhigh, \lvb) {$g$};
\node[downsample, right=0.3cm of gB] (dsBgB) {$\downarrow\!2$};
\draw[arrowline] (gB) -- (dsBgB);
\node[coeff=gdlorange, right=0.45cm of dsBgB] (dJ2) {$d_{J-2,k}$};
\draw[arrowline] (dsBgB) -- (dJ2);
\node[below=2pt of dJ2, font=\scriptsize, text=gdlorange!70!black] {$\in W_{J-2}$};
\def\eqx{8.8}
\node[font=\small, text=gdlpurple!70!black, anchor=west] (eq1) at (\eqx, \lva)
  {$V_{J} = V_{J-1} \oplus W_{J-1}$};
\node[font=\small, text=gdlpurple!70!black, anchor=west] (eq2) at (\eqx, \lvb)
  {$V_{J-1} = V_{J-2} \oplus W_{J-2}$};
\def\arrowx{-3.6}
\draw[gdlgreen!70, thick, ->] (\arrowx, {\lva + 0.5}) -- (\arrowx, {\lvb - 0.5});
\node[font=\scriptsize, text=gdlgreen!70!black, anchor=east] at ({\arrowx - 0.1}, \lva) {fine};
\node[font=\scriptsize, text=gdlgreen!70!black, anchor=east] at ({\arrowx - 0.1}, \lvb) {coarse};
\pgfmathsetmacro{\nesty}{\lvb - 1.2}
\node[font=\small, text=gdlblue!70!black, fill=white, inner sep=3pt,
      rounded corners=2pt, draw=gdlblue!30] at (2.0, \nesty)
  {$\cdots \subset V_{J-2} \subset V_{J-1} \subset V_J \subset \cdots \subset L^2(\R)$};
\end{tikzpicture}
\caption[Wavelet multiresolution decomposition]{Wavelet multiresolution
decomposition: a signal is iteratively decomposed through low-pass ($h$) and
high-pass ($g$) filters followed by subsampling ($\downarrow 2$), separating the
approximation ($c_{j,k}$) and detail ($d_{j,k}$) components at multiple scales.}
\label{fig:grids:wavelet-mra}
\end{figure}

\begin{example}[Haar wavelets]
\label{ex:grids:haar}
The Haar wavelets are the simplest:
\[
\phi(x) = \begin{cases} 1 & 0 \leq x < 1, \\ 0 & \text{otherwise}, \end{cases}
\qquad
\psi(x) = \begin{cases} 1 & 0 \leq x < 1/2, \\ -1 & 1/2 \leq x < 1, \\ 0 & \text{otherwise}. \end{cases}
\]
\end{example}

\begin{example}[Daubechies wavelets]
\label{ex:grids:daubechies}
The Daubechies wavelets $D_N$ have compact support $[0, 2N-1]$ and $N$ vanishing
moments:
\[
  \int_{-\infty}^{\infty} x^k\, \psi(x)\, \mathrm{d}x = 0, \qquad k = 0, 1, \ldots, N-1.
\]
\end{example}

\begin{definition}[Spectral wavelets on graphs]
\label{def:grids:graph-wavelets}
For a graph with Laplacian $L$ and eigenvectors $\{\phi_k\}$, the spectral
wavelets are defined as
\[
  \psi_{s,i} = \sum_{k=0}^{n-1} g(s\lambda_k)\, \phi_k(i)\, \phi_k,
\]
where $g$ is a mother wavelet function and $s$ is the scale parameter. A common
choice is $g(x) = x\, e^{-x}$, which provides both spectral and spatial
localization on the graph.
\end{definition}

Wavelets offer specific advantages for neural networks: multiresolution analysis
that is natural for hierarchical architectures; better localization than Fourier
for local features; sparse representations of natural signals; and approximate
translation invariance depending on the wavelet type.

\begin{definition}[Morlet wavelet]
\label{def:grids:morlet}
The \emph{Morlet wavelet} is a Gaussian-modulated wave:
\[
  \psi_{\omega_0}(t) = c_\sigma\, \pi^{-1/4}\, e^{i\omega_0 t}\, e^{-t^2/2},
\]
where $\omega_0 \geq 5$ is the central frequency and $c_\sigma$ is a
normalization constant. Its Fourier transform is
$\hat{\psi}_{\omega_0}(\omega) = c_\sigma\, \pi^{-1/4}\, e^{-(\omega - \omega_0)^2/2}$,
a Gaussian centred at $\omega_0$.
\end{definition}

\begin{remark}[Optimal time--frequency resolution]
\label{rem:grids:morlet-resolution}
The Morlet wavelet attains equality in the Heisenberg--Gabor uncertainty
principle, $\Delta t \cdot \Delta \omega = \frac{1}{2}$, making it the wavelet
with the best simultaneous resolution in time and frequency. At scale $a$, the
wavelet $\psi_{a,b}(t) = a^{-1/2}\psi((t-b)/a)$ has temporal resolution
$\Delta t \propto a$ and frequency resolution $\Delta \omega \propto 1/a$,
exactly implementing the property of adaptive multiresolution.
\end{remark}

The deepest connection between wavelets and CNNs is established through Mallat's
\emph{scattering transform}~\citep{Mallat2012scattering}, which shows that the
first layers of a CNN can be understood as a bank of wavelet filters with a
modulus nonlinearity.

\begin{definition}[Scattering transform]
\label{def:grids:scattering}
Let $\psi$ be a mother wavelet and $\phi$ a scaling (low-pass) function. The
\emph{scattering operator} of order $m$ is defined recursively:
\begin{align}
S_0 f &= f \star \phi_J, \nonumber \\
S_1 f &= \{\,\abs{f \star \psi_{j_1}} \star \phi_J\,\}_{j_1}, \nonumber \\
S_m f &= \{\,\abs{\cdots\abs{f \star \psi_{j_1}} \star \psi_{j_2}\cdots \star \psi_{j_m}} \star \phi_J\,\}_{j_1 < j_2 < \cdots < j_m},
\end{align}
where $\psi_j(t) = 2^{-j}\psi(2^{-j}t)$ are the dilated wavelets and $\phi_J$ is
the scaling function at the coarsest resolution.
\end{definition}

\begin{proposition}[Properties of the scattering transform]
\label{prop:grids:scattering}
The scattering operator $Sf = (S_0 f, S_1 f, S_2 f, \ldots)$ satisfies:
\begin{enumerate}
  \item \textbf{Translation invariance:} $S[\tau_c f] = S[f]$ for every
    translation $\tau_c f(t) = f(t - c)$.
  \item \textbf{Lipschitz stability to deformations:} for every diffeomorphism
    $\tau$ near the identity,
    $\norm{Sf - S(f \circ \tau)} \leq C\norm{\nabla \tau}_\infty \norm{f}$.
  \item \textbf{Energy conservation:} $\sum_{m=0}^{\infty} \norm{S_m f}^2 = \norm{f}^2$.
\end{enumerate}
\end{proposition}

\begin{proof}
(1) Translation commutes with convolution:
$\tau_c f \star \psi_j = \tau_c(f \star \psi_j)$. The modulus preserves this
translation: $\abs{\tau_c(f \star \psi_j)} = \tau_c\abs{f \star \psi_j}$.
Finally, convolution with $\phi_J$ (low-pass at scale $2^J$) averages over
translations smaller than $2^J$, removing the dependence on $c$.

(2) Since $\psi_j$ has frequency support concentrated at $\abs{\omega} \sim 2^j$,
the convolution $f \star \psi_j$ is sensitive to variations of $f$ at scale
$2^{-j}$. A deformation $\tau$ with $\norm{\nabla \tau}_\infty \leq \epsilon$
shifts frequencies by at most $\epsilon \cdot 2^j$, which is small compared with
the bandwidth $2^j$ of the filter, producing an error proportional to $\epsilon$.

(3) This follows from the Parseval identity for wavelets:
$\norm{f}^2 = \norm{f \star \phi_J}^2 + \sum_j \norm{f \star \psi_j}^2$. Applying
this iteratively to the coefficients $\abs{f \star \psi_j}$ and using that the
modulus preserves the $L^2$ norm yields total energy conservation.
\end{proof}

\begin{remark}[Scattering as a CNN with fixed weights]
\label{rem:grids:scattering-cnn}
The scattering transform architecture is isomorphic to a CNN with the following
correspondences: predefined (not learned) wavelets $\psi_j$ play the role of
convolutional kernels; the modulus $\abs{\cdot}$ plays the role of ReLU;
convolution with $\phi_J$ plays the role of average pooling; and the scattering
order $m$ plays the role of the number of layers. The key difference is that the
scattering transform has theoretical guarantees of stability and invariance,
whereas trained CNNs learn them implicitly.
\end{remark}

\subsection{Spectral convolutions in CNNs}
\label{sec:grids:spectral-cnns}

Spectral convolutions in CNNs exploit the properties of the frequency domain to
implement convolution operations efficiently, especially for large kernels or in
applications with specific geometric structure.

\begin{definition}[Parametrized spectral filter]
\label{def:grids:spectral-filter}
A \emph{parametrized spectral filter} is a function
$h_\theta \colon \R_+ \to \R$ where $\theta \in \R^p$ are trainable parameters.
The filter acts on the frequency spectrum through
\[
  h_\theta(\lambda) = \sum_{k=0}^{K-1} \theta_k\, \phi_k(\lambda),
\]
where $\{\phi_k\}$ is a basis of functions (polynomials, wavelets, etc.).
\end{definition}

\begin{definition}[Spectral convolutional layer]
\label{def:grids:spectral-layer}
A spectral convolutional layer transforms an input signal $x \in \R^n$ through
\[
  y = \sigma(h_\theta \star x + b),
\]
where $h_\theta \star x = \Ft^{-1}[h_\theta(\cdot) \cdot \Ft[x]]$ is the
spectral convolution, $\sigma$ is the activation function, and $b$ is the bias
term.
\end{definition}

\begin{theorem}[Spatial--spectral equivalence]
\label{thm:grids:spatial-spectral}
For domains with regular structure (grids, tori), spatial and spectral
convolutions are mathematically equivalent:
\[
  (f \conv g)(x) = \Ft^{-1}[\Ft[f] \cdot \Ft[g]](x).
\]
They differ, however, in computational complexity (spatial $O(n \cdot k^2)$ vs
spectral $O(n \log n)$), in localization (spatial is local, spectral is global),
and in interpretability (spatial is geometric, spectral is frequency-based).
\end{theorem}

\begin{proof}
The equivalence follows from the convolution theorem for the Fourier transform.
For functions $f, g \in L^1(\R^d)$ (or their discrete analogues on $\Z^d$), the
Fourier transform satisfies $\Ft[f \conv g] = \Ft[f] \cdot \Ft[g]$. Applying
$\Ft^{-1}$ to both sides yields the stated identity. The complexity differences
are: direct spatial convolution performs $O(n \cdot k^2)$ operations (sweeping
$n$ positions applying a kernel of size $k \times k$), while spectral
convolution via the FFT requires $O(n \log n)$ for the forward transform, $O(n)$
for the pointwise multiplication, and $O(n \log n)$ for the inverse transform,
totalling $O(n \log n)$. Localization differs because a finite spatial kernel of
size $k$ corresponds to a spectral filter of full support, and vice versa.
\end{proof}

\paragraph{Practical implementation of spectral CNNs.}
For an input $x \in \R^{H \times W}$ and parametrized filter $h_\theta$, the
spectral forward pass computes $\hat{x} = \mathrm{FFT2D}(x)$, applies the filter
$\hat{y} = h_\theta(\omega) \odot \hat{x}$ (pointwise product), inverts
$y = \mathrm{IFFT2D}(\hat{y})$, and activates $z = \sigma(y + b)$. To avoid edge
effects, circular padding or zero-padding is used; FFT/IFFT may accumulate
floating-point errors; and both spatial and spectral representations must be
stored.

\begin{definition}[Transfer function of a linear CNN]
\label{def:grids:transfer-function}
For a linear CNN (without nonlinear activations) with $L$ layers, the global
transfer function in the frequency domain is
$H(\omega) = \prod_{l=1}^L H_l(\omega)$, where $H_l(\omega)$ is the frequency
response of layer $l$. This factorization is valid only in the absence of
nonlinearities between layers.
\end{definition}

\begin{theorem}[Computational complexity]
\label{thm:grids:complexity}
For a 2D convolution of size $n \times n$ with a $k \times k$ kernel, direct
convolution costs $O(n^2 k^2)$, spectral convolution costs $O(n^2 \log n)$, and
the spectral method is more efficient when $k > O(\log n)$.
\end{theorem}

\begin{proof}
Direct 2D convolution computes, for each of the $n^2$ image positions, a sum of
$k^2$ products, giving $O(n^2 k^2)$ operations. Spectral convolution requires:
(i) 2D FFT of the image, $O(n^2 \log n)$; (ii) 2D FFT of the kernel (with
zero-padding), $O(n^2 \log n)$; (iii) pointwise multiplication in frequency,
$O(n^2)$; (iv) 2D IFFT, $O(n^2 \log n)$. The total is $O(n^2 \log n)$. The
break-even point is obtained by equating $n^2 k^2 = n^2 \log n$, that is,
$k^2 = \log n$, so spectral convolution is preferable when
$k > O(\sqrt{\log n})$.
\end{proof}

\begin{example}[Practical efficiency analysis]
\label{ex:grids:efficiency}
For a $512 \times 512$ image and a $k \times k$ kernel: with $k = 3$, direct
$\sim 2.3 \times 10^6$ ops vs spectral $\sim 2.6 \times 10^6$ ops; with $k = 15$,
direct $\sim 5.9 \times 10^7$ vs spectral $\sim 2.6 \times 10^6$; with $k = 51$,
direct $\sim 6.8 \times 10^8$ vs spectral $\sim 2.6 \times 10^6$. The spectral
cost is essentially constant in $k$.
\end{example}

\begin{theorem}[Spectral universal approximation]
\label{thm:grids:spectral-universal}
Let $\mathcal{G} = (V, E)$ be a finite graph with Laplacian $L$ and spectrum
contained in $[0, \lambda_{\max}]$. For any continuous spectral filter
$h \colon [0, \lambda_{\max}] \to \R$ and $\epsilon > 0$, there exists a
polynomial filter $p(\lambda) = \sum_{k=0}^{K} \theta_k\, T_k(\tilde{\lambda})$
of sufficiently large degree $K$ such that
$\norm{h(\lambda) - p(\lambda)}_{\infty, [0, \lambda_{\max}]} < \epsilon$.
Consequently, the filtered signal satisfies $\norm{h(L)f - p(L)f} \leq
\epsilon \norm{f}$ for every signal $f$ on the graph.
\end{theorem}

\begin{proof}[Sketch]
The spectral interval $[0, \lambda_{\max}]$ is compact, so the result is a
direct consequence of the Weierstrass approximation theorem: every continuous
function on a compact interval can be uniformly approximated by polynomials. The
bound on the filtered signal follows from $\norm{h(L)f - p(L)f} =
\norm{(h-p)(L)f} \leq \norm{h - p}_\infty \norm{f}$, since the eigenvalues of
$(h-p)(L)$ are bounded by $\norm{h-p}_\infty$.
\end{proof}

The main limitations of spectral convolutions are their non-locality (spectral
filters affect the signal globally), the requirement of domains with FFT
structure (regular grids), difficulties with non-periodic boundary conditions,
and reduced interpretability compared with spatial convolution.

\subsection{Frequency interpretation of classical CNNs}
\label{sec:grids:freq-interpretation}

The frequency analysis of traditional CNNs reveals deep insights into what
features they learn and how they process visual information.

\begin{definition}[Frequency response of a kernel]
\label{def:grids:freq-response}
For a convolutional kernel $K \in \R^{k \times k}$, its frequency response is
\[
  H(\omega_x, \omega_y) = \sum_{i=0}^{k-1} \sum_{j=0}^{k-1} K[i,j]\,
    e^{-2\pi i(\omega_x i + \omega_y j)}.
\]
\end{definition}

\begin{example}[Analysis of common kernels]
\label{ex:grids:common-kernels}
The Sobel edge detector $K_{\text{Sobel}} = \begin{psmallmatrix} -1 & 0 & 1 \\
-2 & 0 & 2 \\ -1 & 0 & 1 \end{psmallmatrix}$ has a high-pass frequency response
with a maximum at medium horizontal frequencies. The Gaussian filter
$K_{\text{Gauss}}[i,j] = \frac{1}{2\pi\sigma^2}
e^{-\frac{(i-c)^2 + (j-c)^2}{2\sigma^2}}$ has a low-pass response with a smooth
cutoff determined by $\sigma$. The Laplacian
$K_{\text{Lap}} = \begin{psmallmatrix} 0 & -1 & 0 \\ -1 & 4 & -1 \\ 0 & -1 & 0
\end{psmallmatrix}$ has an isotropic high-pass response.
\end{example}

\begin{remark}[Spectral evolution during training]
\label{rem:grids:spectral-evolution}
Empirical observations show that, during the training of a CNN, the spectral
distribution of the filters evolves following a characteristic pattern: early
layers learn low-pass filters and edge detectors (high frequencies);
intermediate layers learn band-pass filters for specific textures; and deep
layers learn complex filters that are less spectrally interpretable. This
behaviour has been documented experimentally across multiple architectures,
although it lacks a general formal proof. Related empirical studies show that
CNNs exhibit a bias towards low frequencies: they learn low-frequency components
first and require more iterations to capture high-frequency detail.
\end{remark}

The spectral methods of this section find direct applications in domains where
the frequency structure of the data is fundamental: medical image reconstruction
(MRI, computed tomography), scientific signal processing (astronomy,
seismology), and neural-network compression through spectral pruning. The focus
of this book, however, is on the mathematical foundation of these methods.

\subsection{Limitations and synthesis}
\label{sec:grids:spectral-limits}

The classical spectral convolutions of this section require a regular structure
(uniform grids for an efficient FFT), periodic boundary conditions or adequate
padding, and a fixed size --- consequently restricting them to regular Euclidean
domains. They also exhibit a trade-off between localization and efficiency, as
summarized in \Cref{tab:grids:methods}: spatial convolution offers excellent
localization but only moderate efficiency, while spectral convolution offers
excellent efficiency at the cost of localization; polynomial approximations and
wavelets occupy intermediate positions.

\begin{table}[h]
\centering
\caption{Comparison of convolutional methods.}
\label{tab:grids:methods}
\begin{tabular}{@{}lccc@{}}
\toprule
\textbf{Method} & \textbf{Localization} & \textbf{Efficiency} & \textbf{Interpretability} \\
\midrule
Spatial convolution     & Excellent & Medium    & High   \\
Spectral convolution    & Poor      & Excellent & Medium \\
Polynomial approximation & Good     & Good      & Medium \\
Wavelets                & Excellent & Medium    & High   \\
\bottomrule
\end{tabular}
\end{table}

Further limitations of spectral methods are numerical instability (accumulation
of errors over multiple FFT/IFFT, aliasing in subsampling without anti-aliasing
filters, floating-point precision especially on GPUs), reduced interpretability
(spectral filters are less visually intuitive and their relation to geometric
features is indirect), and the absence of spectral transferability between
domains: a filter $h(\lambda)$ learned on the spectrum
$\{\lambda_k^{(1)}\}$ of one domain has no canonical interpretation on the
spectrum $\{\lambda_k^{(2)}\}$ of another unless a natural correspondence between
eigenvalues exists. This last limitation motivates the use of polynomial
approximations such as Chebyshev filters, which define the filter as a function
of the operator $h(L)$ rather than acting on individual eigenvalues.

\begin{proposition}[Fundamental localization--globality trade-off]
\label{prop:grids:loc-global}
Let $h(\lambda)$ be a spectral filter supported on
$[\lambda_{\min}, \lambda_{\max}]$ and let
$K(x,y) = \sum_k h(\lambda_k)\, \phi_k(x)\,\phi_k(y)$ be its associated spatial
kernel. Then
\[
  \supp_{\mathrm{freq}}(h) \cdot \supp_{\mathrm{spatial}}(K) \geq C > 0,
\]
where $C$ depends only on the dimension of the domain. In particular, a filter
perfectly localized in frequency ($h = \delta_{\lambda_k}$) is completely
delocalized in space ($K(x,y) = \phi_k(x)\phi_k(y)$, global support).
\end{proposition}

\begin{proof}
This is a direct consequence of the uncertainty principle for the spectral
transform: the concentration of a function and of its spectral representation
cannot both be made arbitrarily small. For the case $h = \delta_{\lambda_k}$,
the kernel reduces to $K(x,y) = \phi_k(x)\phi_k(y)$, which is non-zero almost
everywhere on the domain (since the eigenvectors of the Laplacian vanish only on
a set of measure zero). The general case follows by a duality argument analogous
to the classical Heisenberg principle.
\end{proof}

Harmonic analysis thus provides a unifying framework that connects classical
CNNs (Fourier analysis on regular grids), spectral graph networks (spectral
analysis on graphs), geometric CNNs (harmonic analysis on groups and manifolds),
and transformers (attention as adaptive convolution). This perspective reveals
that many seemingly different architectures share deep mathematical foundations.

These examples make the chapter's thesis concrete. Convolutional networks are
not an ad hoc engineering trick but the canonical realization of a single
principle --- respect the translation symmetry of the grid --- read in two dual
languages: the spatial language of weight sharing and the spectral language of
the Fourier basis. The spatial description is local and geometric; the spectral
description is the one that \emph{survives the loss of the grid}, since on a
graph there is no notion of sliding a kernel, yet there is a Laplacian whose
eigenvectors generalize the Fourier basis. Holding the principle fixed and
enlarging the symmetry group --- or abandoning the grid altogether --- is
precisely the programme of the chapters that follow, beginning in
\Cref{ch:groups} with networks equivariant to rotations and reflections.
